%-------------------------
% Cover Letter - Pearson Senior Software Engineer
% Author: Ajay Venkatesh
%-------------------------

\documentclass[letterpaper,11pt]{article}

\usepackage[top=0.8in, bottom=0.8in, left=0.9in, right=0.9in]{geometry}
\usepackage[hidelinks]{hyperref}
\usepackage{lmodern}
\usepackage{parskip}
\usepackage{microtype}

\setlength{\parskip}{6pt}
\setlength{\parindent}{0pt}

\begin{document}

\begin{flushleft}
\textbf{\large Ajay Venkatesh} \\
Brooklyn, NY \\
+1 516 585 9013 \\
\href{mailto:av3855@nyu.edu}{av3855@nyu.edu}
\end{flushleft}

\vspace{10pt}

May 15, 2026

\vspace{6pt}

Hiring Manager \\
Pearson

\vspace{10pt}

Dear Hiring Manager,

My name is Ajay Venkatesh --- finishing my M.S. in Computer Engineering at NYU, with several years of production software experience: a few years at PwC in Bengaluru, a summer SDE internship at Amazon, and ongoing on-campus work at NYU's Novel AI Technologies. My focus has been on \textbf{Java backend services, data-intensive applications,} and the patient end-to-end ownership that public-facing software demands --- which is the same triangle the Senior Software Engineer role at Pearson lives in.

The posting reads like a description of how I already work. At Amazon I engineered backend services in \textbf{Java} (Coral, Smithy, Dagger, CBOR) with low-level design patterns and unit/integration testing via \textbf{Mockito} --- well-documented APIs, code-review discipline, and quality treated as part of the build, not an afterthought. I provisioned the supporting infrastructure on \textbf{AWS CDK} (Lambda, API Gateway, DynamoDB, IAM, CloudWatch) with multi-stage CI/CD and rollback controls, and built event-driven services on \textbf{ECS Fargate, SQS, and SNS} that materially reduced response latency. The same posture maps cleanly to \textbf{API-first, service-oriented} systems serving statewide programs at Pearson's scale.

On data-intensive systems specifically: my Big Data project built a \textbf{PySpark / HDFS} pipeline over millions of flight records, training Random Forest and Gradient Boosted Tree models with rigorous EDA --- feature engineering at scale, anomaly detection, and pulling signal out of messy real-world data. At NYU I engineered \textbf{ETL pipelines} transforming \textbf{NoSQL} into \textbf{PostgreSQL} with concurrency control, processing hundreds of thousands of records, and powered \textbf{AWS QuickSight} dashboards downstream. At PwC I spent three years on a distributed backend on \textbf{Microsoft Azure} with \textbf{SQL Server}, parallelizing batch ingestion that cut execution from hours to minutes. Diagnosing bottlenecks in data-heavy systems is where I find the work most rewarding.

On frontend: I've shipped \textbf{React} interfaces at Amazon (React 18 with optimized API integration and client-side caching) and at NYU (analytics platform with role-based access for paying clients). \textbf{Vue.js} would be a new framework for me, but the underlying patterns --- component composition, reactive state, build tooling --- are the same shape, and the JD framing of "comparable modern frontend framework with willingness to learn Vue" is exactly the kind of bridge I'd want to make.

What pulls me toward Pearson is the mission and the seriousness of the domain. Education software runs on real consequences --- statewide contracts, public-facing programs, students whose progress depends on the system being correct under load. That's an environment where engineering judgment compounds, and where data-intensive, performance-aware work actually matters in real units.

I'd welcome the chance to discuss further. Thanks for considering my application.

\vspace{12pt}

Sincerely, \\
Ajay Venkatesh

\end{document}
